\FloatBarrier
\section{Discussion}
The positive-measure formulation makes several distinctions operational. A discretization defines how mass is represented; regularization defines which admissible reconstructions are preferred; optimization determines how accurately a chosen criterion is solved; and identification concerns what the data support about the generating system. Their diagnostics differ. Increasing a display grid to 256 bins can reduce quantization error without providing 256 independent diffusion degrees of freedom. Reducing a residual can improve prediction while worsening distributional shape or changing the apparent number of species.

The support selector addresses a narrower failure: diffuse low-amplitude assignment of a shared atomic rate across unrelated frequencies. Residualization asks whether a proposed rate contributes beyond the other candidate decays, and aggregation can retain weak but coherent evidence. The unrestricted candidate remains available, so a hard screen is not imposed regardless of its predictive consequences. Automatic order is chosen from observations, with a null model and an explicit cap. Neither the true synthetic order nor a visually preferred component height enters selection.

There are important limits to this construction. The groups are geometric frequency windows, not chemical resonances obtained from a validated deconvolution. Screening a whole window assumes aggregate positive evidence; an overly large window can dilute a localized signal. Estimated nuisance rates can leave systematic residuals, and the screen does not account for model uncertainty. At near collision, residualization becomes poorly conditioned. Retaining support in that case avoids interpreting numerical collapse as absence, but cannot supply the missing information.

The prediction allowance is also a decision rule rather than a confidence procedure. Proposition~\ref{prop:paired} concerns a fixed pair of predictions independent of validation noise. Choosing the best candidate, reusing data in the mask, and searching multiple supports invalidate a direct inferential reading of one standard deviation. Reporting an independent external test protects the empirical comparison from that particular reuse, but does not calibrate the selected order. Independent mask construction, nested resampling, uncertainty over supports, and tests under correlated or estimated noise are reasonable subsequent studies, not properties established here.

The present evidence does not isolate every mechanism. Native-to-S comparisons change support restriction and final order selection together. Native baselines fit the 24 outer training rows directly, whereas the S procedure uses 18 rows for candidate construction followed by a 24-row refit. Both receive the same outer training information, but this is not a matched computation-budget ablation. A factorial comparison of screening only, predictive selection only, and the combined procedure would be required to apportion their effects. Since both native paths often reach the same final estimate, the study does not establish a special benefit of DOME's moment exchanges or RAI's particular starts. A global optimum is not certified, and a successful candidate search is not an exhaustive search over all continuous supports.

MF and atomic inversion answer different reconstruction questions. A positive factor profile can capture a broad polymer-like distribution, while a discrete model estimates a finite union of decay rates. Simplex normalization prevents the trivial scaling ambiguity but does not identify the factors chemically. The archived MF rank threshold counts resolved matrix directions; using that count as a cap can exclude weaker directions. This is particularly relevant when component spectra are nearly proportional. An improved automatic species selector cannot be obtained merely by relabeling a factor count.

The transfer illustration is useful because it reveals rather than hides this mismatch. RAI-S and DOME-S draw narrow bands for finite-width truth because their models are atomic. Their reduced spectral leakage is a benefit in a particular endpoint, not proof that the widths have been recovered. Similarly, lower binned transport loss for a policy within a limited domain does not demonstrate general neural superiority. The very small independent-test gain after RAI-Net retraining provides a concrete negative result.

The main empirical limitation is external validity. All new tests are synthetic, use known iid Gaussian noise, and operate after a signal mask. Baseline errors, phase errors, gradient calibration, relaxation contamination, exchange, convection, nonuniform spectral noise, and solvent suppression can violate this contract. No new experimental DOSY validation, molecular identity assignment, or molecular-weight calibration is included. The archived multi-method collection contains other implementations and adaptations, but its heterogeneous configurations are not substituted for a matched external comparator study here. Reproducibility of the current simulations and competence against strong independently tuned controls are different requirements.

\section{Conclusions}
Joint positive Laplace inversion benefits from keeping its mathematical objects explicit: masses, densities, atomic supports, and shared factors are not interchangeable. The derived results establish limited statements about discretization, uniqueness, and conditional selection statistics, while the implementation makes model order and spectral support data-dependent within a finite search. In the frozen atomic test, the shared support procedure reduces leakage and clean prediction error for both proposal paths, with modest improvement in full recovery and a small increase in unmatched mass. Proportional and weak mixtures remain difficult, and distributional controls expose the cost of an atomic assumption. These findings support an auditable selection framework and a bounded empirical improvement; they do not establish universal component identification.
